\section{The Documented Patient Concerns}

\subsection{Where the twenty-one come from}

Two dated research passes on July 28, 2026 supply the evidence base for this
section. The first, by Gemini 3.1 Pro, returned six concern families. The second,
by ChatGPT 5.6 Thinking Extended, returned sixteen numbered concerns, each with
the question a patient is likely to ask and what the trial team should explain,
together with thirteen supporting references. The two overlap on five items.
De-duplicated, they give twenty-one distinct documented concerns, and those
twenty-one are the subject of the rest of this paper.

Neither pass invented its material. Behind them sits human evidence. A 2026
qualitative study of 50 previously operated patients found that 56 percent would
consider semi-autonomous robot-assisted surgery, and that willingness depended
strongly on which parts of the procedure the robot controlled; the principal
themes were informed consent, additional risks, surgeon-patient trust, and
anticipated benefit \cite{WuSemiAI2026}. A cross-sectional survey of 330
oncology patients found approximately half expressing concern about artificial
intelligence in cancer care, led by loss of human interaction and by medical
error \cite{TelesAI2026A}. A systematic review of patient and public
perspectives on robotic surgery found that expectation and experience diverge
systematically, and that expectation is usually set by promotion rather than by
evidence \cite{Jauniaux2025}. A survey of public perceptions found that the most
common single misconception is that the robot operates by itself
\cite{BrarRAS2024X}.

Two conclusions follow, and this paper is built on both. First, patients are not
opposed to the technology; they are opposed to being asked to accept it without
specifics. Second, the specifics they want are not general. They are per-step,
per-limit, per-name.

\vspace{0.30em}

\begin{xltabular}{\textwidth}{L{4.2cm} C{1.9cm} Y}
\toprule
\textbf{Concern family} & \textbf{Concerns} & \textbf{Where this paper answers it} \\
\midrule
\endfirsthead
\multicolumn{3}{@{}l}{\footnotesize\itshape Table continued from the previous page.}\\
\toprule
\textbf{Concern family} & \textbf{Concerns} & \textbf{Where this paper answers it} \\
\midrule
\endhead
\midrule
\multicolumn{3}{r@{}}{\footnotesize\itshape Table continued on the next page.}\\
\endfoot
\bottomrule
\endlastfoot
Control and the human element & 5 & \S7.3, \S7.4; Figures 8, 18, 19 \\
Safety and technical failure & 5 & \S7.4, \S7.5, \S7.6; Figures 20, 21 \\
Evidence and effectiveness & 4 & \S4.2, \S5.1, \S10; Figures 13, 27, 28 \\
Data, privacy, and security & 3 & \S9.3, \S9.4; Figures 17, 25 \\
Fairness and applicability & 2 & \S6.3, \S8.3; Figures 14, 23 \\
Accountability and burden & 2 & \S11, \S12; Figures 29, 30 \\
\end{xltabular}

\vspace{0.30em}

\begin{pafloat}
\begin{pafig}[d2-type: nested containers]
% ---- family 1 ----
\node[d2key,text width=26mm] (c1) at (0,0)      {1. Who is driving\\\S7.3, Fig 19};
\node[d2key,text width=26mm] (c2) at (0,-1.5)   {2. Can it be stopped\\\S7.4, Fig 18};
\node[d2box,text width=26mm] (c3) at (0,-3.0)   {3. Will the surgeon defer\\\S7.3, Fig 19};
\node[d2box,text width=26mm] (c4) at (0,-4.5)   {4. Will I have a doctor\\\S3.1, Fig 8};
\node[d2box,text width=26mm] (c5) at (0,-6.0)   {5. Less time with my team\\\S6.2, Fig 15};
% ---- family 2 ----
\node[d2key,text width=26mm] (s1) at (5.4,0)    {6. Unintended motion\\\S7.4, Fig 21};
\node[d2box,text width=26mm] (s2) at (5.4,-1.5) {7. Misreading a margin\\\S7.5, Fig 21};
\node[d2key,text width=26mm] (s3) at (5.4,-3.0) {8. Vessel injury\\\S7.5, Fig 20};
\node[d2box,text width=26mm] (s4) at (5.4,-4.5) {9. Instrument failure\\\S7.6, Fig 21};
\node[d2box,text width=26mm] (s5) at (5.4,-6.0) {10. Conversion to open\\\S5.3, Fig 12};
% ---- family 3 ----
\node[d2gray2,text width=26mm](e1) at (10.8,0)   {11. Cancer control\\\S10.1, Fig 27};
\node[d2box,text width=26mm] (e2) at (10.8,-1.5) {12. Among the first\\\S5.1, Fig 13};
\node[d2gray2,text width=26mm](e3) at (10.8,-3.0){13. Newer as better\\\S10.4, Fig 28};
\node[d2box,text width=26mm] (e4) at (10.8,-4.5) {14. Team experience\\\S11.2, Fig 29};
% ---- family 4 ----
\node[d2key,text width=26mm] (d1) at (0,-9.4)   {15. Who sees the video\\\S9.3, Fig 25};
\node[d2box,text width=26mm] (d2) at (0,-10.9)  {16. Training other models\\\S9.3, Fig 25};
\node[d2key,text width=26mm] (d3) at (0,-12.4)  {17. Could it be attacked\\\S9.4, Fig 17};
% ---- family 5 ----
\node[d2gray2,text width=26mm](f1) at (5.4,-9.4) {18. People like me\\\S6.3, Fig 14};
\node[d2key,text width=26mm] (f2) at (5.4,-10.9) {19. Software change\\\S8.3, Fig 23};
% ---- family 6 ----
\node[d2gray2,text width=26mm](b1) at (10.8,-9.4) {20. Who is accountable\\\S11.1, Fig 29};
\node[d2gray2,text width=26mm](b2) at (10.8,-10.9){21. What will it cost\\\S12.2, Fig 30};
\begin{scope}[on background layer]
\node[d2cont,fit=(c1)(c2)(c3)(c4)(c5)] (F1) {};
\node[d2cont,fit=(s1)(s2)(s3)(s4)(s5)] (F2) {};
\node[d2cont,fit=(e1)(e2)(e3)(e4)]     (F3) {};
\node[d2cont,fit=(d1)(d2)(d3)]         (F4) {};
\node[d2cont,fit=(f1)(f2)]             (F5) {};
\node[d2cont,fit=(b1)(b2)]             (F6) {};
\end{scope}
\node[d2title,anchor=south west] at (F1.north west) {concerns.control\_and\_the\_human\_element};
\node[d2title,anchor=south west] at (F2.north west) {concerns.safety\_and\_technical\_failure};
\node[d2title,anchor=south west] at (F3.north west) {concerns.evidence\_and\_effectiveness};
\node[d2title,anchor=south west] at (F4.north west) {concerns.data\_privacy\_and\_security};
\node[d2title,anchor=south west] at (F5.north west) {concerns.fairness\_and\_applicability};
\node[d2title,anchor=south west] at (F6.north west) {concerns.accountability\_and\_burden};
\legkey{-1.6}{-14.2}{protoblue}{answered by a hard numeric limit, 7 of 21}
\legkey{4.2}{-14.2}{protowhite}{answered by a procedural guarantee, 9 of 21}
\legkey{10.0}{-14.2}{pagraym}{governance or disclosure only, 5 of 21}
\end{pafig}
\vspace{-0.7cm}
\figcaption{Figure 5. The twenty-one documented concerns grouped into six disjoint families, with\\
the answering section and figure printed inside each item, and a fill key stating\\
how many are answered by a limit, a guarantee, or by governance alone.}
\end{pafloat}

\subsection{Control and the human element}

Five concerns sit here, and they are the most frequently raised of the
twenty-one. Who is actually driving the robot. Whether the surgeon can stop it,
and how fast. Whether the surgeon will defer to the model even when their own
judgement differs. Whether the participant will still have a doctor who knows
them. Whether relying on the system will reduce time with the care team.

The first is answered by classification and by architecture. The system is a
Class II collaborative device operating under continuous human oversight, and
the on-premises model has no path to an actuator: it emits a candidate plan, a
deterministic gate tests that plan, and the operating surgeon signs or declines
it. Section 7.3 sets this out and Figure 19 draws every message. The public
misconception that the robot operates by itself \cite{BrarRAS2024X} is not
corrected by asserting the opposite; it is corrected by showing a diagram in
which the arrow that would be needed is absent.

The third concern, deference, is the subtlest and the least often addressed. It
is not enough that the surgeon may decline; the question is whether they will.
This protocol records every override with its reason, publishes the override
rate as an exploratory endpoint, and requires no justification for declining, so
that declining carries no procedural cost. Automation bias is a measured
quantity here rather than an assumed absence \cite{FDATrans2024}.

The fourth and fifth concerns are about relationship rather than technology.
Every route into the study notifies the participant's own treating oncologist,
on every branch, including the branch where the participant turns out to be
ineligible. The direct sponsor channel of \S6.2 is additional, never a
replacement, and Figure 15 draws the notification on each path.

\subsection{Safety and technical failure}

Five concerns: unintended motion, misreading tissue or a margin, injury to a
major vessel, instrument failure during the procedure, and conversion to open
surgery.

The first three are answered by limits that are enforced below the level of the
software making the decision. Per-arm tip force is capped at 3 N and the sum
across arms at 18 N. Positional tolerance is 2 mm and force reproducibility is
0.5 N. A vascular no-fly corridor extends 4 mm either side of the superior
mesenteric and portal veins, and a plan that would enter it is rejected by a
deterministic gate before the surgeon ever sees it as executable. These are the
kinds of limits the international standard for robotically assisted surgical
equipment contemplates \cite{iec8060127}, and the fail-safe design standard for
autonomous and semi-autonomous systems frames the stopping behaviour
\cite{ieee7009}.

Instrument failure and conversion are answered differently, because they are not
prevented, they are planned for. An open tray and backup instruments are in the
room for every procedure. Conversion to open surgery is a recorded outcome, not
a failure of the study or of the participant, it does not end participation, and
the decision belongs to the operating surgeon alone. Section 7.7 gives the
comparator context.

\subsection{Evidence and effectiveness}

Four concerns: whether the operation removes the cancer as completely as
standard surgery, whether the participant is among the first humans to receive
it, whether newer is being presented as better, and how experienced the team is.

This is the family this paper answers least completely, and \S10 says so with a
figure rather than a hedge. Figure 7 places cancer-control effectiveness in the
quadrant of concerns that are raised often and answered only by governance or
disclosure. The reason is structural: eighteen participants cannot establish an
oncologic outcome. What the protocol can do, and does, is report R0 rate, ISGPS
grade B/C fistula rate, and 90-day mortality against a contemporaneous
nationwide comparator of 1000 robotic pancreatoduodenectomies
\cite{DutchCohort2025}, and restrict enrollment to a site already past the third
phase of that learning curve, so the curve is not paid for by the participant.

On the third concern, the paper takes the position the surveyed literature
recommends: use neutral language and present absolute evidence rather than
adjectives \cite{Jauniaux2025}. Nowhere in this document is the system described
as revolutionary, state of the art, or cutting edge. The FDA's own information
page states that cancer-related long-term outcomes have not been established for
every use of robotically assisted surgical devices \cite{FDARobot2022}, and that
statement is reproduced rather than argued with.

\subsection{Data, privacy, and security}

Three concerns: who sees the recording of the operation, whether the data will
train other systems, and whether the system can be attacked.

The answers are architectural and are drawn rather than asserted. Section 9.3
publishes a read list per store and a retention period per store, in years.
Section 9.4 states the network position: during a procedure there is no route
out of the building, so the attack surface is the building. Three things cross
the boundary before a procedure, a signed model build, a protocol version, and
the pinned registry entry; two cross afterwards, reports to the FDA and the IRB;
and nothing crosses during. Figure 17 draws that boundary and Figure 25 draws
the pipeline. The relevant FDA guidance on cybersecurity in medical devices
\cite{FDACyber2026} and the joint transparency principles for machine
learning-enabled devices \cite{FDATrans2024} are the frame.

Data is not used to train a commercial model without a separate consent. That is
a statement about something that does not happen, and Figure 25 renders it as an
absent edge, which is the only honest way to draw a negative.

\subsection{Fairness, applicability, accountability, and burden}

Four concerns close the set: whether the system was tested on people like the
participant, whether the software can change after consent, who is accountable
if something goes wrong, and what participation will cost.

On applicability the honest answer is uncomfortable and is given in \S6.3: at
$n \leq 18$ subgroup performance is not estimable, and what the protocol commits
to is reporting the composition of the enrolled population by age, sex, race,
ethnicity, body mass index, and prior therapy, consistent with FDA guidance on
reporting such data in device studies \cite{FDADiver2017}, together with the
real-world performance monitoring the agency has sought comment on
\cite{FDARealW2025}.

On software change the answer is a hard one: the version is frozen at consent
and recorded in a registry, and any change forces a re-consent. Until the
participant acts, no study procedure proceeds under the changed version.
Figure 23 draws consent as the state machine this implies.

Accountability and cost are the two concerns that land in the honest-gap
quadrant alongside effectiveness. Section 11 answers accountability with a rule,
one accountable party per failure class, and Figure 29 publishes the matrix.
Section 12 answers cost line by line and marks two items as unsettled: lost
income, and continued drug access after study closure \cite{NCICosts2024}.

\begin{pafloat}
\begin{pafig}[graphviz-type: graph, bipartite, full page]
\foreach \i/\lab in {1/{1. Who is driving},2/{2. Can it be stopped},3/{3. Will it be deferred to},4/{4. Will I have a doctor},5/{5. Less time with my team},6/{6. Unintended motion},7/{7. Misreading a margin},8/{8. Vessel injury},9/{9. Instrument failure},10/{10. Conversion to open},11/{11. Cancer control},12/{12. Among the first},13/{13. Newer as better},14/{14. Team experience},15/{15. Who sees the video},16/{16. Training other models},17/{17. Could it be attacked},18/{18. People like me},19/{19. Software change},20/{20. Who is accountable},21/{21. What will it cost}}{
  \node[gvbox,text width=30mm,minimum height=5mm,inner sep=1.8pt] (k\i) at (0,{-(\i-1)*0.72}) {\lab};}
\node[gvboxk,text width=44mm,minimum height=5mm,inner sep=1.8pt] (a1)  at (8.8,-0.05) {Class II collaborative device, continuous oversight, no model-to-actuator path};
\node[gvboxk,text width=44mm,minimum height=5mm,inner sep=1.8pt] (a2)  at (8.8,-0.90) {Emergency stop: $\leq 3$ ms cross-arm, $\leq 500$ ms system-wide};
\node[gvbox,text width=44mm,minimum height=5mm,inner sep=1.8pt]  (a3)  at (8.8,-1.72) {Override logged with reason; declining needs no justification};
\node[gvbox,text width=44mm,minimum height=5mm,inner sep=1.8pt]  (a4)  at (8.8,-2.54) {Treating oncologist notified on every branch of every route};
\node[gvboxk,text width=44mm,minimum height=5mm,inner sep=1.8pt] (a5)  at (8.8,-3.42) {Tip force $\leq 3$ N per arm, $\leq 18$ N summed across arms};
\node[gvboxk,text width=44mm,minimum height=5mm,inner sep=1.8pt] (a6)  at (8.8,-4.30) {Vascular no-fly gating, SMV and portal vein, 4 mm corridor};
\node[gvboxk,text width=44mm,minimum height=5mm,inner sep=1.8pt] (a7)  at (8.8,-5.18) {Positional tolerance $\leq 2$ mm, force reproducibility $\leq 0.5$ N};
\node[gvbox,text width=44mm,minimum height=5mm,inner sep=1.8pt]  (a8)  at (8.8,-6.06) {Conversion is a recorded outcome, not a failure; backup tray in room};
\node[gvbox,text width=44mm,minimum height=5mm,inner sep=1.8pt]  (a9)  at (8.8,-6.94) {R0 rate, ISGPS B/C rate, and 90-day mortality against a nationwide comparator};
\node[gvboxk,text width=44mm,minimum height=5mm,inner sep=1.8pt] (a10) at (8.8,-7.88) {Phase 0 gate: $\geq 1000$ simulations, $\geq 2$ frameworks, USL $\geq 7.0$};
\node[gvbox,text width=44mm,minimum height=5mm,inner sep=1.8pt]  (a11) at (8.8,-8.82) {Sentinel stagger: no second participant at a new dose until the first window closes};
\node[gvbox,text width=44mm,minimum height=5mm,inner sep=1.8pt]  (a12) at (8.8,-9.76) {Site restricted past learning-curve phase 3; operator counts published};
\node[gvbox,text width=44mm,minimum height=5mm,inner sep=1.8pt]  (a13) at (8.8,-10.64){Read list published per store; retention stated in years};
\node[gvbox,text width=44mm,minimum height=5mm,inner sep=1.8pt]  (a14) at (8.8,-11.46){No commercial model training without a separate consent};
\node[gvboxk,text width=44mm,minimum height=5mm,inner sep=1.8pt] (a15) at (8.8,-12.28){No network route out of the building during a procedure};
\node[gvbox,fill=pagraym,text width=44mm,minimum height=5mm,inner sep=1.8pt] (a16) at (8.8,-13.10) {Subgroup composition reported per FDA guidance; exclusions named};
\node[gvbox,text width=44mm,minimum height=5mm,inner sep=1.8pt]  (a17) at (8.8,-13.92){Version frozen at consent; any change forces re-consent};
\node[gvbox,fill=pagraym,text width=44mm,minimum height=5mm,inner sep=1.8pt] (a18) at (8.8,-14.68) {One accountable party per failure class, named in \S11};
\node[gvbox,fill=pagraym,text width=44mm,minimum height=5mm,inner sep=1.8pt] (a19) at (8.8,-15.44) {Sponsor pays study-only items; two items remain unsettled};
\begin{scope}[on background layer]
\node[gvcluster,fit=(k1)(k21),inner sep=4pt]  (KK) {};
\node[gvcluster2,fit=(a1)(a19),inner sep=4pt] (AA) {};
\end{scope}
\node[gvctitle,anchor=south west] at (KK.north west) {the twenty-one documented concerns};
\node[gvctitle,anchor=south west] at (AA.north west) {the clause, limit, or number that answers it};
\foreach \s/\t in {1/a1,2/a2,3/a3,3/a1,4/a4,5/a4,6/a5,6/a2,7/a7,8/a6,9/a8,10/a8,11/a9,12/a10,12/a11,13/a9,14/a12,15/a13,16/a14,17/a15,19/a17}{
  \draw[gvundirb] (k\s.east) -- (\t.west);}
\foreach \s/\t in {18/a16,20/a18,21/a19}{
  \draw[gvundird] (k\s.east) -- (\t.west);}
\node[font=\tiny\sffamily\itshape,text=protogray,anchor=north west,align=left]
  at (-1.9,-15.9) {Solid blue: answered by a hard numeric limit or an architectural fact. Dashed gray: answered by governance or disclosure alone.\\No edge joins two concerns and no edge joins two clauses, so no concern in this figure is answered with another concern.};
\end{pafig}
\vspace{-0.7cm}
\figcaption{Figure 6. Every documented concern wired to the clause, limit, or number that answers\\
it, drawn as a strictly bipartite graph so that no concern is answered with another\\
concern, and dashed where the answer is governance or disclosure alone.}
\end{pafloat}

\subsection{How completely each concern is answered}

Figure 7 places sixteen of the documented concerns on two independent axes. The
horizontal axis is prevalence, scored from how often the concern appears across
the four surveyed human sources. The vertical axis is answer completeness,
scored 1.0 where the protocol supplies a numeric limit, 0.7 where it supplies a
procedural guarantee, 0.5 where it supplies a governance commitment, and 0.3
where it supplies disclosure only.

Four concerns land in the lower-right quadrant: raised often, answered softly.
They are cancer-control effectiveness, accountability for harm, loss of the
human relationship, and cost. This paper does not argue them out of that
quadrant. Sections 10, 11, and 12 state, for each, exactly what would have to be
true for the answer to become a hard one, and what this study can contribute
towards making it so.

\begin{pafloat}
\begin{pafig}[mermaid-type: quadrantChart]
\fill[pablue2]  (5.6,5.6) rectangle (11.2,11.2);
\fill[protowhite] (0,5.6) rectangle (5.6,11.2);
\fill[pagrayl]  (0,0) rectangle (5.6,5.6);
\fill[pagraym]  (5.6,0) rectangle (11.2,5.6);
\draw[protogray,line width=0.4pt] (0,0) rectangle (11.2,11.2);
\draw[protogray,line width=0.5pt] (5.6,0) -- (5.6,11.2);
\draw[protogray,line width=0.5pt] (0,5.6) -- (11.2,5.6);
\node[font=\tiny\sffamily\bfseries,text=protoblue,anchor=north east,align=right] at (11.0,11.0) {asked often, answered hard\\lead with these};
\node[font=\tiny\sffamily\bfseries,text=protoblue,anchor=north west,align=left]  at (0.2,11.0)  {answered hard, asked less\\keep available};
\node[font=\tiny\sffamily\bfseries,text=protogray,anchor=south west,align=left]  at (0.2,0.2)   {softly answered, asked less\\monitor};
\node[font=\tiny\sffamily\bfseries,text=protoblue,anchor=south east,align=right] at (11.0,0.2)  {asked often, answered softly\\the honest gap};
\node[font=\scriptsize\sffamily,anchor=north] at (5.6,-0.5) {prevalence in the surveyed sources $\longrightarrow$};
\node[font=\scriptsize\sffamily,rotate=90,anchor=south] at (-0.6,5.6) {completeness of the protocol answer $\longrightarrow$};
\foreach \x/\y/\lx/\ly/\anc/\lab in {%
  10.42/9.86/11.30/9.10/west/{who is driving the robot},
  8.96/10.75/8.20/11.05/east/{override and rescue speed},
  9.86/10.30/10.90/10.55/west/{malfunction, unintended motion},
  6.94/10.08/6.20/10.55/east/{vascular injury and bleeding},
  10.08/6.50/10.95/6.90/west/{cancer-control effectiveness},
  7.95/6.94/7.10/7.30/east/{surgeon and team experience},
  9.41/6.16/10.30/5.70/west/{who is accountable for harm},
  7.39/9.63/6.55/9.95/east/{recording and secondary use},
  5.49/10.53/4.60/10.85/east/{cybersecurity and network reach},
  6.50/4.93/5.60/5.30/east/{algorithmic bias and my anatomy},
  3.81/9.18/2.95/9.55/east/{software change after consent},
  4.59/6.72/3.70/7.05/east/{automation bias in the surgeon},
  8.51/7.39/9.40/7.85/west/{being among the first humans},
  9.74/4.26/10.60/3.80/west/{loss of the human relationship},
  6.16/4.03/5.30/3.60/east/{hype rather than evidence},
  8.18/4.59/7.30/4.20/east/{cost and post-trial burden}}{
  \draw[pagrayd,line width=0.3pt] (\x,\y) -- (\lx,\ly);
  \fill[protoblue] (\x,\y) circle (0.085);
  \node[font=\tiny\sffamily,anchor=\anc,inner sep=1.6pt,fill=protowhite,fill opacity=0.85,text opacity=1] at (\lx,\ly) {\lab};}
\node[font=\scriptsize\sffamily\itshape,text=protogray,anchor=north west,align=left]
  at (0,-1.2) {Scoring: 1.0 a numeric limit, 0.7 a procedural guarantee, 0.5 a governance commitment, 0.3 disclosure only. Four concerns\\sit in the lower-right quadrant. Sections 10, 11, and 12 state what would have to change for each of the four to move.};
\end{pafig}
\vspace{-0.7cm}
\figcaption{Figure 7. Sixteen documented concerns positioned by how often they are raised in the\\
surveyed literature against how completely this protocol answers them, with the lower-right\\
quadrant collecting the concerns answered by governance or disclosure alone.}
\end{pafloat}

\subsection{The five answered by governance alone, named}

Five of the twenty-one are answered by governance or disclosure and by nothing
harder. Naming them is more useful than distributing them quietly through the
text, because a reader who knows which five they are can weigh them
deliberately.

\vspace{0.30em}

\begin{xltabular}{\textwidth}{C{0.7cm} L{4.2cm} Y}
\toprule
\textbf{\#} & \textbf{Concern} & \textbf{Why the answer cannot be harder than governance} \\
\midrule
\endfirsthead
\multicolumn{3}{@{}l}{\footnotesize\itshape Table continued from the previous page.}\\
\toprule
\textbf{\#} & \textbf{Concern} & \textbf{Why the answer cannot be harder than governance} \\
\midrule
\endhead
\midrule
\multicolumn{3}{r@{}}{\footnotesize\itshape Table continued on the next page.}\\
\endfoot
\bottomrule
\endlastfoot
12 & Will this work better than the standard operation & A single-arm study of at most 18 participants cannot answer a comparative question. The honest answer is a comparator table and a stated inability, \S 10.5 \\
13 & Long-term cancer outcomes & Progression-free and overall survival are followed to 24 months and reported descriptively. No inference about durability is available at this size, \S 10.4 \\
18 & People like me & The eligibility criteria exclude the frailer patients in whom the question is sharpest. Subgroup performance is not estimable at n at most 18, and \S 7.3 says so \\
20 & Who is accountable & Accountability is assigned by a published responsibility matrix and by named parties, not by a numeric limit. It is enforceable, but it is governance, \S 11 \\
21 & What it will cost & Study-only items are paid by the sponsor. Two items, lost income and post-closure data access, remain unsettled and are printed unsettled, \S 12.3 \\
\end{xltabular}

\vspace{0.30em}

\noindent The distinction being drawn is not between strong answers and weak
ones. A governance answer with a named accountable party, a published matrix,
and an escalation route is a real answer. The distinction is that a numeric
limit fails visibly and a governance commitment fails quietly, so the second
kind requires the participant to know where to look. That is what \S 11.6 is
for.

Two of the five could become harder answers, and it is worth being specific
about how. Concern 12 becomes answerable only in a randomised study, which this
is not and which cannot be run until feasibility is established; that is the
argument for this study existing. Concern 21 becomes answerable the moment the
two unsettled cost items are settled, which is a decision the sponsor can take
and has not yet taken. The other three are structural at this phase.

\subsection{The same safeguards, read as capabilities you hold}

Every safeguard in the parent protocol is written with the sponsor as the actor:
the sponsor verifies, the sponsor reports, the sponsor halts. Figure 8 redraws
the same twelve safeguards as a use case diagram with the participant as the
primary actor. Nothing in the protocol changed. Only the actor did, and the
result reads differently: ten of the twelve are capabilities the participant
exercises.

The two that are not are pressing the emergency stop and issuing a software
version notice. Both require someone physically present in the room or inside
the build system, and neither can be exercised against the participant: the stop
can only halt, and the notice can only trigger a re-consent that the participant
is free to decline.

\begin{pafloat}
\begin{pafig}[plantuml-type: use case, @startuml]
\umlactor{pt}{-7.4}{-4.6}{\textbf{Participant}\\(primary actor)}
\node[umlusekey] (u1)  at (-2.4,0)    {Consent, or decline with no consequence};
\node[umlusekey] (u2)  at (-2.4,-2.3) {Opt out of the Physical AI component};
\node[umlusekey] (u3)  at (-2.4,-4.6) {Set the autonomy level accepted per step};
\node[umlusecase](u4)  at (-2.4,-6.9) {Inspect the Phase 0 simulation result};
\node[umlusekey] (u5)  at (2.6,0)     {Require a named human to approve every motion};
\node[umlusecase](u6)  at (2.6,-2.3)  {Trigger a stop through the surgeon};
\node[umlusekey] (u7)  at (2.6,-4.6)  {Accept or decline the daraxonrasib restart};
\node[umlusecase](u8)  at (2.6,-6.9)  {Read my own audit trail};
\node[umlusekey] (u9)  at (7.6,0)     {Withdraw at any visit, no reason required};
\node[umlusecase](u10) at (7.6,-2.3)  {Choose the fate of banked specimens};
\node[umlusecase](u11) at (7.6,-4.6)  {Be told when the software version changes};
\node[umlusecase](u12) at (7.6,-6.9)  {Receive my day-30 and day-90 results};
\begin{scope}[on background layer]
\node[umlpkg,fit=(u1)(u4)(u9)(u12)] (SYS) {};
\end{scope}
\node[umlpkgtab,anchor=south west] at (SYS.north west) {Phase 1 PDAC Physical AI trial};
\foreach \u in {u1,u2,u3,u4,u5,u7,u9,u10,u11,u12}{\draw[umlassoc] (pt) -- (\u);}
\umlactor{sg}{12.6}{-0.6}{Operating surgeon}
\umlactor{ds}{12.6}{-3.6}{DSMB}
\umlactor{si}{12.6}{-6.6}{Sponsor-Investigator}
\draw[umlassoc] (sg) to[out=180,in=0,looseness=0.85] (u5);
\draw[umlassoc] (sg) to[out=190,in=10,looseness=0.85] (u6);
\draw[umlassoc] (ds) to[out=180,in=-10,looseness=0.85] (u6);
\draw[umlassoc] (si) -- (u11);
\draw[umlassoc] (si) to[out=190,in=0,looseness=0.85] (u12);
\draw[umlinherit] (u6) -- node[umlguard,pos=0.5]{$\ll$extend$\gg$} (u5);
\draw[umldash] (u8) -- node[umlguard,pos=0.5]{$\ll$include$\gg$} (u4);
\draw[umldash] (u11) to[out=150,in=-30,looseness=0.9] node[umlguard,pos=0.5]{$\ll$include$\gg$ forces re-consent} (u1);
\node[umlnote,text width=64mm,anchor=north west] at (-8.6,-9.2)
  {Ten of the twelve capabilities are exercised by the participant. The two that are not, pressing the stop and issuing a version notice, are the two that require someone in the room or inside the build system. Neither can be exercised against the participant.};
\node[umlnote,text width=52mm,anchor=north west] at (1.6,-9.2)
  {In the parent protocol these twelve appear as sponsor obligations. Nothing in the protocol changed here. Only the actor did.};
\end{pafig}
\vspace{-0.7cm}
\figcaption{Figure 8. The twelve safeguards of the protocol redrawn as a use case diagram with\\
the participant as the primary actor, so that each one reads as a capability they hold\\
rather than as an obligation the sponsor has undertaken on their behalf.}
\end{pafloat}

\subsection{Where each concern physically lives}

A concern with no physical address is unanswerable. A concern that resolves to a
named cabinet, a named console, or a named cable can be inspected, and
inspection is what the surveyed patients asked for. Figure 9 assigns thirteen of
the twenty-one to the component that answers them, grouped into the operating
room, the on-premises cabinet, the governance bodies, and everything outside the
building.

The figure's central claim is again an absence. There is a dashed line between
the on-premises model and the network boundary, and it is struck through. During
a procedure no route exists out of the building. That single struck line answers
concern 17 more completely than any paragraph in this section.

\begin{pafloat}
\begin{pafig}[diagrams (python)-type: Cluster + Node]
\dgnodew{dgtiled}{-6.6}{0}{\glyphrobot}{Eight robotic arms, force-capped}{n1}
\dgnode{dgtile}{-3.8}{0}{\glyphsignal}{Sensor stack, 640 channels}{n2}
\dgnodew{dgtiled}{-6.6}{-2.4}{\glyphmon}{Surgeon console, sole authority}{n3}
\dgnode{dgtile}{-3.8}{-2.4}{\glyphdoc}{Procedure recorder}{n4}
\dgnode{dgtile}{0.8}{0}{\glyphai}{On-premises model, advisory}{n5}
\dgnodew{dgtiled}{3.6}{0}{\glyphshield}{Deterministic safety gate}{n6}
\dgnode{dgtile}{0.8}{-2.4}{\glyphdb}{Version registry}{n7}
\dgnode{dgtile}{3.6}{-2.4}{\glyphlock}{Hash-chained audit trail}{n8}
\dgnode{dgtile}{2.2}{-4.8}{\glyphdb}{Data vault}{n9}
\dgnodeg{dgtileg}{8.0}{0}{\glyphdoc}{DSMB}{n10}
\dgnodeg{dgtileg}{8.0}{-2.4}{\glyphdoc}{IRB, open to you}{n11}
\dgnodeg{dgtileg}{8.0}{-4.8}{\glyphuser}{Sponsor-Investigator}{n12}
\dgnodeg{dgtilek}{2.2}{-8.0}{\glyphnet}{Network boundary, closed}{n13}
\dgnodeg{dgtilek}{6.0}{-8.0}{\glyphcloud}{Hospital network and cloud}{n14}
\begin{scope}[on background layer]
\node[dgcluster,fit=(n1)(n1l)(n2)(n2l)(n3)(n3l)(n4)(n4l)] (C1) {};
\node[dgcluster,fit=(n5)(n5l)(n6)(n6l)(n7)(n7l)(n8)(n8l)(n9)(n9l)] (C2) {};
\node[dgcluster2,fit=(n10)(n10l)(n11)(n11l)(n12)(n12l)] (C3) {};
\node[dgcluster2,fit=(n13)(n13l)(n14)(n14l)] (C4) {};
\end{scope}
\node[dgctitle,anchor=south west] at (C1.north west) {operating room, sterile field};
\node[dgctitle,anchor=south west] at (C2.north west) {on-premises cabinet, same building};
\node[dgctitle2,anchor=south west] at (C3.north west) {governance, off the critical path};
\node[dgctitle2,anchor=south west] at (C4.north west) {everything outside the building};
\foreach \n/\c in {n1/{6, 9},n2/{7},n3/{1, 2, 3, 10},n4/{15},n6/{8},n7/{19},n8/{20},n9/{16},n13/{17}}{
  \node[circle,fill=pagraym,draw=protoblack,line width=0.3pt,inner sep=0.6pt,
        font=\tiny\sffamily,anchor=south west] at ([xshift=1mm,yshift=1mm]\n.north east) {\c};}
\draw[dgedgeb] (n3) -- node[font=\tiny\sffamily,above]{approved motion only} (n1);
\draw[dgedge]  (n2) -- (n3);
\draw[dgedge]  (n1) to[out=-90,in=180,looseness=0.85] (n4);
\draw[dgedged] (n5) -- (n6);
\draw[dgedgeb] (n6) to[out=180,in=20,looseness=0.9] node[font=\tiny\sffamily,above,pos=0.6]{passed plans only} (n3);
\draw[dgedge]  (n7) -- (n5);
\draw[dgedge]  (n6) -- (n8);
\draw[dgedge]  (n4) to[out=-90,in=180,looseness=0.85] (n9);
\draw[dgedged] (n8) to[out=0,in=180,looseness=0.85] (n10);
\draw[dgedged] (n9) to[out=0,in=200,looseness=0.85] (n12);
\draw[dgedged] (n9) to[out=20,in=190,looseness=0.85] (n11);
\draw[dgedged] (n13) -- node[font=\tiny\sffamily,above]{no route} (n14);
\draw[dgedgek] (n5) to[out=-100,in=100,looseness=0.9] (n13);
\pxmark{1.34}{-5.6}
\node[font=\tiny\sffamily,fill=protowhite,inner sep=1.5pt,anchor=east] at (1.1,-5.6) {no route exists during a procedure};
\end{pafig}
\vspace{-0.7cm}
\figcaption{Figure 9. Each documented concern resolved to the physical component that answers it,\\
grouped into the operating room, the on-premises cabinet, the governance bodies, and\\
everything outside, with the intra-procedural route drawn and struck through.}
\end{pafloat}

\subsection{The twenty-one, one row each}

The table below is the paper's index. It states each concern in the form
participants state it, the question they are most likely to ask, the answer this
protocol gives, and the class of that answer. A reader who wants only one thing
from \S 3 should take this table.

\vspace{0.30em}

\begin{xltabular}{\textwidth}{C{0.7cm} L{3.4cm} L{4.2cm} Y}
\toprule
\textbf{\#} & \textbf{Concern} & \textbf{The question asked} & \textbf{The answer, and its class} \\
\midrule
\endfirsthead
\toprule
\textbf{\#} & \textbf{Concern} & \textbf{The question asked} & \textbf{The answer, and its class} \\
\midrule
\endhead
\bottomrule
\endfoot
\bottomrule
\endlastfoot
1 & Who is driving the robot & Is my surgeon controlling every movement, or approving a plan, or supervising an automated task? & Approving. Every motion is surgeon-signed; the model has no actuator path. Hard architectural limit. \\
2 & Can the surgeon stop it, and how fast & Can the surgeon stop the system immediately, and how quickly is human control restored? & At most 3 ms across arms, at most 500 ms system-wide, measured on every procedure. Hard numeric limit. \\
3 & Will the surgeon defer to the model & Will my surgeon follow the model even when their own judgement differs? & Overrides are logged with a reason and the override rate is a published endpoint. Procedural guarantee. \\
4 & Will I still have a doctor who knows me & Will I still have a surgeon who knows me and is personally responsible? & The operating surgeon is named in your consent and is accountable for every intra-operative motion. Procedural guarantee. \\
5 & Less time with my care team & Will relying on the system reduce my face-to-face time? & The direct sponsor channel is additional; every branch notifies your own oncologist. Procedural guarantee. \\
6 & Malfunction or unintended motion & Could the system move unexpectedly and injure something? & Tip force capped at 3 N per arm and 18 N summed, enforced in hardware. Hard numeric limit. \\
7 & Misreading tissue or a margin & Could it misidentify a structure or a tumour boundary? & Positional tolerance 2 mm, force reproducibility 0.5 N, both gate-enforced. Hard numeric limit. \\
8 & Injury to a major vessel & Could an autonomous movement cause severe bleeding? & A 4 mm no-fly corridor around the SMV and portal vein; a plan entering it is rejected before it is executable. Hard numeric limit. \\
9 & Instrument failure mid-procedure & What if an instrument breaks while I am open? & Backup instruments and an open tray in the room for every procedure; conversion rehearsed in the checklist. Procedural guarantee. \\
10 & Conversion to open surgery & What happens if the robotic approach has to be abandoned? & A recorded outcome, not a failure; participation continues; you are told before discharge. Procedural guarantee. \\
11 & Cancer control & Will this remove my cancer as completely as standard surgery? & R0 rate, fistula rate, and 90-day mortality reported against a 1000-case nationwide comparator. Comparator, not a result. \\
12 & Being among the first & Am I one of the first people to receive this? & Phase 0 gate of 1000 simulations across two frameworks precedes any patient; sentinel stagger precedes any second dose. Hard gate. \\
13 & Newer described as better & Is this actually better, or just newer? & No superlative appears in this paper; the FDA's own caution is reproduced verbatim. Disclosure. \\
14 & Team experience & How many of these has my surgeon done, on this exact platform? & Case count, platform count, autonomy level, and trainee participation published before consent. Procedural guarantee. \\
15 & Who sees the recording & Who will see the video of my operation? & A published read list per store; operating surgeon, Site PI, and you. Procedural guarantee. \\
16 & Training other models & Will my data be used to train other systems? & Not without a separate consent; drawn in Figure 25 as an absent edge. Procedural guarantee. \\
17 & Could it be attacked & Could someone interfere with the system during my operation? & No network route out of the building during a procedure. Hard architectural limit. \\
18 & People like me & Was it tested on people with my anatomy, age, and history? & Composition reported per FDA guidance; subgroup performance not estimable at $n \leq 18$. Disclosure. \\
19 & Software change after consent & Could the software change after I agree? & Version frozen at consent; any change forces a re-consent before any procedure. Procedural guarantee. \\
20 & Who is accountable & Who answers if I am harmed? & Exactly one accountable party per failure class, named in \S 11 and published before consent. Governance. \\
21 & What will it cost & What will I pay, during and after? & Line by line in \S 12; two items, lost income and post-closure drug access, remain unsettled. Governance. \\
\end{xltabular}

\vspace{0.30em}

\noindent Seven of the twenty-one are answered by a hard numeric or
architectural limit, nine by a procedural guarantee, and five by governance or
disclosure alone. Those proportions are the honest shape of the answer, and they
are the same proportions Figure 5 colours and Figure 7 plots.
